\title[Intrinsic reconstruction of webs]{Intrinsic reconstruction of webs from nonreduced multiplication-failure schemes}
\author{Qian Qi}
\date{September 23, 2026 \quad (Revision 135)}
\subjclass[2020]{14M12, 14J28, 13C40, 14B05, 14D20, 05E05}
\keywords{Intrinsic reconstruction, webs of quadrics, Fitting schemes, Schur functors, exterior support, Jacobian quartics, nilpotent filtration}
\hypersetup{pdftitle={Intrinsic reconstruction of webs from nonreduced multiplication-failure schemes, revision 135},pdfauthor={Qian Qi}}
\begin{abstract}
We prove that the abstract nonreduced multiplication-failure scheme
of a cube-zero algebra with Hilbert function \((1,4,6)\) determines
its relation web up to projective equivalence whenever the web is
basepoint-free and its Jacobian quartic is smooth.  No further
genericity condition is required.  The intrinsic deepest normal cone
recovers both irreducible Pluecker components of the web, transported
by one projective transformation.  A contraction-kernel theorem
computes the exterior support of the Jacobian component: a smooth
quartic has support ten, whereas secant and tangent four-vectors have
support at most eight.  This excludes all exceptional recombinations;
indeed three essential variables suffice.  The support of the ambient rank-defect locus is contained in the
inverse image of the binary-quartic subspace variety; no equality
or ambient component classification is asserted.

The first nilpotent layer recovers the polarized Jacobian K3.  The
full nonreduced scheme separates the web classes not separated by
that polarized surface.  A separate technical supplement retains the full primary-boundary
and relative-specialization theory.  The reconstruction proof uses
the nonreduced normal cone directly, without a higher-corank primary
atlas or an exceptional-open hypothesis.
\end{abstract}
\maketitle
